\subsection{Square polarization removes equality dummies}
Suppose a logical value currently has exactly two factor aliases $a,b$,
and its equality penalty uses a fresh dummy $d$:
\[
 C(a,b,d)=a^2+b^2+d^2-ab-bd-da.
\]
The reason for retaining this dummy can be a prescribed square $a^2$
inside the incidence graph: directly substituting the two-alias penalty
$Q(a,b)$ may collide with its degree-four diagonal contribution. The new
operation first changes selected prescribed squares into products of
same-logical aliases, $a^2\mapsto ab$ (or $b^2\mapsto ab$). It then replaces
$C(a,b,d)$ by $Q(a,b)$ and deletes the dummy coordinate $d$.

\begin{theorem}\label{thm:polarization}
Make any finite selection of such same-logical prescribed-square changes,
and replace each affected two-alias equality cycle by $Q$. If the full
collected final polynomial has coefficient height one, it is a height-one
quartic with an integer zero set bijective to the original arithmetic
source. Each deleted equality dummy reduces its dimension by one.
\end{theorem}
\begin{proof}
Changing a prescribed edge's label from $a^2$ to $ab$ leaves its two
incidence endpoints and orientation unchanged. The graph energy therefore
remains exactly $\frac12\sum_v d_v^2$, with its divergences evaluated using
the modified labels. It is still nonnegative, and all labels still have
degree at most two. Also
\[
 Q(a,b)=(a-b)^2(a^2+ab+b^2)\ge0,
\]
and $Q(a,b)=0$ if and only if $a=b$: the second factor equals
$(a+b/2)^2+3b^2/4$ and vanishes only at $a=b=0$.

At any zero of the final energy, every nonnegative equality term is zero,
so same-logical aliases agree and units equal one. Consequently every
polarized label $ab$ has exactly the value previously specified by $a^2$.
The occurrence-boundary identity is unchanged after logical substitution.
Vanishing vertex divergences therefore give all normalized source rows,
and monic reconstruction gives an actual arithmetic source zero.

Conversely, for an original source zero assign both $a$ and $b$ their
common logical value. All polarized edge labels then have their old values,
all new $Q$ terms vanish, and the same unique subtree flow assignment makes
every vertex divergence zero. Deleted dummies have unique old value $d=a=b$,
so extension and deletion are inverse on the zero sets. This also proves
that the source-polynomial lift identity \eqref{eq:networklift} survives
universally, not only at source zeros.

The degree is at most four by construction. Exact degree four and
coefficient height one are finite properties of the entire collected
polynomial and are explicitly checked. They do not follow merely from the
simplicity of the incidence graph after monomial substitution. The operation
is accepted only after this full collection check, including all equality
and unit terms. Thus possible collisions do not invalidate the claim.
\end{proof}
The theorem makes no assertion that an arbitrary polarization is safe for
height one, nor that a chosen finite sequence is optimal. It separates the
unconditional witness proof from the coefficient test selecting a particular
literal endpoint.

\subsection{Selected outputs and matched source controls}
\begin{center}\small
\begin{tabular}{@{}lrrrr@{}}\toprule
Endpoint & Variables & Terms & Height & Constant\\\midrule
\file{minimum528} &528&20816&1&1\\
\file{compact534} &534&18543&1&1\\
\file{balanced580} &580&12001&1&1\\
\file{sparse620} &620&8318&1&1\\
\file{sparser717} &717&7176&1&1\\\bottomrule
\end{tabular}
\end{center}
All five have exact degree four, unrestricted integer coordinates and every
declared coordinate present. Their complete coefficient arrays and literal
polynomial text are supplied. The first two use the specifically identified
\file{exposed247} source with 247 variables, 179 rows, 4889 terms and
height 170496. This is the stage after the first three affine projections,
but before the later height-reducing residual basis changes. It is not the
literal \file{balanced247} equation.

The fixed selection of 24 affine square rows expands this source to
271 variables and 203 rows; coefficient normalization gives 345 variables
and 277 rows. The two final graphs use 20 and 24 hub vertices. Polarization
then removes ten and eleven equality dummies, respectively. These are
arithmetic-source and compiler improvements together. The 534-variable,
18543-term point improves both dimension and support over the Round~18
571-variable, 19833-term point. The 528-variable point saves a further six
coordinates but has more support.

To isolate polarization alone, the replay also rebuilds the exact four
Round~18 graph certificates before making the new change. It checks their
historical hashes and holds the arithmetic source and graph parameters fixed:
\begin{center}\small
\begin{tabular}{@{}lrr@{}}\toprule
Matched source & Before variables/terms & After variables/terms\\\midrule
Round 18 compact &571/19833&559/19809\\
Round 18 balanced &593/12027&580/12001\\
Round 18 sparse &646/8368&620/8318\\
Round 18 sparser &742/7226&717/7176\\\bottomrule
\end{tabular}
\end{center}
Thus each matched-source transformation improves both displayed measures.
The 559-variable control is reproducible but is not selected as a principal
endpoint because the new-source 534-variable point has fewer variables and
fewer terms. The selected 580-variable output retains the old 263-variable
arithmetic source; the last two retain the old 285-variable sparse source.
No new refutation calibration is needed for any of these choices.

